% Generated draft: requires mathematical and source review.
\documentclass[11pt]{article}
\usepackage[margin=25mm]{geometry}
\usepackage{fontspec}
\setmainfont{Latin Modern Roman}
\usepackage{amsmath,amssymb,mathtools,unicode-math}
\setmathfont{Latin Modern Math}
\usepackage{enumitem,needspace,xurl,hyperref}
\setlength{\parindent}{0pt}
\setlength{\parskip}{5pt}
\setlength{\emergencystretch}{3em}
\newcommand{\N}{\mathbb{N}}
\newcommand{\Z}{\mathbb{Z}}
\newcommand{\Q}{\mathbb{Q}}
\newcommand{\R}{\mathbb{R}}
\newcommand{\C}{\mathbb{C}}
\newcommand{\F}{\mathbb{F}}
\newcommand{\A}{\mathbb{A}}
\newcommand{\PP}{\mathbb{P}}
\newcommand{\E}{\mathbb{E}}
\newcommand{\eps}{\varepsilon}
\newcommand{\dd}{\,\mathrm{d}}
\begin{document}
\section*{UnsolvedMath ID: 11000003}
\textbf{Source problem code:} catalogue entry
\subsection{Definitions and mathematical statement}
Let \(\gamma_1,\ldots,\gamma_{2g}\) be the standard basis of the free abelian group \(\Z^{2g}\), and retain the source's space \(Y_g\) and functions \(L_\gamma\). Define
\[
 L=\sum_{i=1}^{2g}L_{\gamma_i}:Y_g\longrightarrow[0,\infty).
\]
The source excerpt does not define \(Y_g\), \(L_\gamma\), or the genus range. The research summary calls \(Y_g\) Torelli space. In the usual closed oriented surface interpretation this would mean Teichm\"uller space modulo the Torelli group, the kernel of the mapping-class-group action on integral first homology, but that setup must be checked against the source.

\emph{Problem.} Study \(L\) as a Morse function: find its critical sets and use them to deduce properties of the topology of \(Y_g\). In particular, give upper and lower estimates for its homotopy dimension. In the usual CW-type convention,
\[
 \operatorname{hdim}(Y_g)
 =\inf\{\dim X:X\text{ is a CW complex homotopy equivalent to }Y_g\}.
\]
The extracted phrase ``homotopy dimension 2 of \(Y_g\)'' contains a detached \(2\), plausibly a footnote marker. It is not treated as a demand to prove that \(Y_g\) has dimension \(2\); its role requires source confirmation.

A rigorous Morse-theoretic formulation must establish the smooth or stratified setting, regularity of \(L\), and the relevant nondegeneracy and global hypotheses. The request for ``critical sets'' does not by itself specify ordinary Morse, Morse--Bott, or another version of Morse theory. Without the definitions of the summands, no differential or explicit critical-point condition can be supplied.
\subsection{Short English statement}
Analyze the sum of the source's \(2g\) basis-indexed functions on \(Y_g\) by Morse theory, determine its critical sets, and use them to bound the homotopy dimension. The space and summands are omitted, and a detached \(2\) may be a footnote marker.
\subsection{Sources}
\begin{itemize}
\item Local UnsolvedMath registry, permanent ID 11000003, source code catalogue entry, Problem 2.2, PDF page 13. No individual proposer is supplied.
\item Farb, \emph{Problems on Mapping Class Groups and Related Topics}, Problem 2.2 at the specified page. Supplied URL: \url{https://www.math.uchicago.edu/~farb/papers/mcgbook.pdf}; recorded access 2026-07-29; extraction pdf-font-structure. The edited volume has 329 explicitly typeset Problem, Conjecture, and Question headings across chapters, with repeated local numbers disambiguated by PDF page.
\item The research summary attributes the program to Farb's chapter (2006), identifies the space as Torelli space, mentions later topology and homological-dimension work without citations, and says no dedicated published Morse treatment of this particular \(L\) was located. These are supplied assessments.
\item Supplied status open and OPEN-TRIAGE. Current primary-literature status and redistribution terms for the publicly accessible author/publisher copy remain NEEDS\_REVIEW.
\end{itemize}
\end{document}
